\documentstyle[%


righttag%


,11pt%


,amssymb%


,russian%               remove in Eng.


,izvru%                 Set izven% in Eng.


]{amsart}





% 





\newcommand{\dist}{\operatorname{dist}}


\newcommand{\tts}[1]{}





\theoremstyle{definition}


\theorembodyfont{\rm}


\newtheorem{notenonum}{\hskip\parindent Замечание}


\renewcommand{\thenotenonum}{}





%\hoffset=-15mm%                  For Eng.





\setcounter{page}{53}


\god{1998}


\nomer{5~(432)} %                        Remove in Eng.


%\nomer{Vol.\,42, No.\,5}%               Set in Eng.


%\str{\pageref{begin}--\pageref{end}}%  Set in Eng.


\udk{517.518}


\author{\it А.-Р.\,К.\,Рамазанов}


% NB:   ^^ remove in Eng.


\title{Знакочувствительные аппроксимации\\


ограниченных функций полиномами}





\date{}


% \pagestyle{myheadings}%         Set in Eng.





\pagestyle{plain} %              Remove in Eng.


\begin{document}





% \thispagestyle{plain}%          Set in Eng.





\maketitle


\label{begin}





Вопросы существования, единственности и устойчивости элемента наилучшего


равномерного приближения для знакочувствительных аппроксимаций изучены в


работах \cite{1}, \cite{2} (см. также \cite{3}--\cite{6}).


Ниже дается решение задачи Е.П.\,Долженко о критериях сколь угодно


точного приближения полиномами ограниченной функции в равномерной


метрике с ограниченным знакочувствительным весом (а также с


полунепрерывным сверху и, в частности, с непрерывным весом); уточняется


теорема об отделении полунепрерывных функций непрерывными (\cite{7},


с.\,390), точнее, дается оценка модуля непрерывности функции, отделяющей


две заданные полунепрерывные функции, причем приводимое нами


доказательство самой теоремы об отделении существенно короче; полученная


оценка применяется для выяснения скорости полиномиального приближения


ограниченной функции в равномерной метрике со знакочувствительным весом.





Знакочувствительным весом на отрезке $\Delta=[a,b]$ называется


упорядоченная пара $p=(p_-,p_+)$


определенных на $\Delta$ неотрицательных функций $p_-(x)$ и $p_+(x)$.


Вес $p=(p_-,p_+)$


называется ограниченным на $\Delta$, если функции $p_\mp(x)$


ограничены на $\Delta$.





Для ограниченных на $\Delta$ функции $f(x)$ и веса $p=(p_-,p_+)$ будем


придерживаться следующих обозначений:


\begin{align*}


\|f\|_\Delta&=\sup\{|f(x)|:x\in\Delta\},\\


\|p\|_\Delta&=\max\{\|p_-\|_\Delta,\;\|p_+\|_\Delta\},\\


|(f,p)(x)|&=f^+(x)p_+(x)+f^-(x)p_-(x),\\


f^+(x)&=\max\{f(x),\;0\},\quad f^-(x)=(-f(x))^+,\\


|f|_{p,\Delta}&=\sup\{|(f,p)(x)|:x\in\Delta\},\\


\Delta(f\ge c)&=\{x\in\Delta:f(x)\ge c\}


\end{align*}


($c$ --- число; аналогично определяются $\Delta(f<c)$, $\Delta(f>c)$).





При $\varepsilon>0$ пусть $F_\varepsilon$ --- пересечение множеств


$\Delta(p_-\ge\varepsilon)$ и $\Delta(p_+\ge\varepsilon)$, $\Delta_-$ и


$\Delta_+$


--- их замыкания соответственно, $\Delta_\varepsilon$ --- пересечение


этих замыканий.


Для $[\alpha,\beta]\subset\Delta$ положим


\begin{align*}


M(\alpha,\beta)&=\sup\{f(x):x\in[\alpha,\beta]\cap\Delta


(p_+\ge\varepsilon)\},\\


m(\alpha,\beta)&=\inf\{f(x):x\in[\alpha,\beta]\cap\Delta


(p_-\ge\varepsilon)\}


\end{align*}


и определим функции


\begin{align*}


M_+(x)&=M_+(x,\varepsilon)=\lim_{\delta\to+0}


M(x-\delta,x+\delta),\quad x\in\Delta_+;\\


m_-(x)&=m_-(x,\varepsilon)=\lim_{\delta\to+0}


m(x-\delta,x+\delta),\quad x\in\Delta_-.


\end{align*}


Так как множества $\Delta(p_-\ge\varepsilon)$ и


$\Delta(p_+\ge\varepsilon)$ плотны в множестве $\Delta_\varepsilon$,


легко видеть, что функция $M_+(x)$ полунепрерывна сверху, а функция


$m_-(x)$ полунепрерывна


снизу на $\Delta_\varepsilon$.





Для ограниченных на отрезке $\Delta$ функции $f(x)$ и веса


$p=(p_-,p_+)$ имеет место





\begin{thm}


Для существования алгебраических полиномов $Q_n(x)$ таких, что


$|f-Q_n|_{p,\Delta}\to 0$ при


$n\to\infty$, необходимо и достаточно выполнения при всех достаточно


малых $\varepsilon>0$ неравенства


\begin{equation}


M_+(x,\varepsilon)\le m_-(x,\varepsilon),\quad


x\in\Delta_\varepsilon.


\end{equation}


\end{thm}





\begin{notenonum}


В случае веса $p=(p_-,p_+)$ с полунепрерывными сверху (в частности,


непрерывными)


на $\Delta$ функциями $p_-(x)$ и $p_+(x)$ при всех $\varepsilon>0$


множества $\Delta(p_\mp\ge\varepsilon)$ замкнуты,


$\Delta_\varepsilon=F_\varepsilon$, причем


$f(x)=M_+(x,\varepsilon)=m_-(x,\varepsilon)$, если


$M_+(x,\varepsilon)\le m_-(x,\varepsilon)$.





Поэтому, как легко видеть из приводимого ниже доказательства


теоремы 1, если на $\Delta$ заданы ограниченная функция $f(x)$ и вес


$p=(p_-,p_+)$ с


ограниченными и полунепрерывными сверху функциями $p_\mp(x)$, то


справедлива


теорема 1 с заменой в ней условия (1) на непрерывность $f(x)$ на


$F_\varepsilon$.


\end{notenonum}





Пусть $\cal P_n$ --- множество всех алгебраических полиномов степени не


выше $n$ ($n=0,1,\dotsc$), и для ограниченных на $\Delta$ функции


$f(x)$ и веса $p=(p_-,p_+)$ положим


$$


E_n(f,p,\Delta)=\inf\{|f-Q_n|_{p,\Delta}:Q_n\in\cal P_n\}.


$$


Вопрос о скорости стремления к нулю $E_n(f,p,\Delta)$ при $n\to\infty$


сводится к следующей задаче.





Пусть на отрезке $\Delta$ для ограниченной полунепрерывной сверху


функции $u(x)$ и


полунепрерывной снизу функции $v(x)$ выполняется неравенство


$u(x)\le v(x)$. Оценить


модуль непрерывности непрерывной функции $\varphi(x)$, для которой


$u(x)\le\varphi(x)\le v(x)$, $x\in\Delta$.





Из известного доказательства (\cite{7}, с.\,390) теоремы о


существовании функции $\varphi(x)$ нельзя получить подобной оценки,


поэтому приведем решение этой задачи.





При $\delta\in[0,\,b-a]$ положим


$$


\Omega(\delta,u,v,\Delta)=\sup\{[u(x)-v(y)]^+:


|x-y|\le\delta,\; x,y\in\Delta\}.


$$


Легко показать, что $\Omega(\delta)=\Omega(\delta,u,v,\Delta)$ не


убывает и полунепрерывна сверху на $[0,\,b-a]$, $\Omega(\delta)$


непрерывна в нуле, $\Omega(0)=0$. На отрезке $[0,\,b-a]$ определим


непрерывную функцию $\Omega_1(\delta)$,


положив $\Omega_1(\delta)=\Omega(\delta)$ при $\delta=0,\,b-a$;


$\Omega_1((b-a)/n)=\sup\{\Omega(\delta):\delta\in((b-a)/n,\;


(b-a)/(n-1))\}$ при


$n=2,3,\dotsc$; $\Omega_1(\delta)$ линейна на отрезках вида


$[(b-a)/n,\;(b-a)/(n-1)]$ ($n=2,3,\dotsc$).


Возьмем теперь модуль непрерывности $\omega(\Omega_1,\delta)$ функции


$\Omega_1(\delta)$ в равномерной метрике и при $x\in\Delta$ рассмотрим


функции


\begin{align*}


\varphi_u(x)&=\sup\{u(y)-\omega(\Omega_1,\,|x-y|):y\in\Delta\},\\


\varphi_v(x)&=\inf\{v(y)+\omega(\Omega_1,\,|x-y|):y\in\Delta\}.


\end{align*}





\begin{thm}


Если на отрезке $\Delta$ для ограниченных полунепрерывной сверху функции


$u(x)$ и


полунепрерывной снизу функции $v(x)$ выполняется неравенство


$u(x)\le v(x)$, то для $\varphi(x)$


$(x\in\Delta)$, равной любой из функций $\varphi_u(x)$ и


$\varphi_v(x)$, имеем


\begin{gather*}


u(x)\le\varphi(x)\le v(x),\\


|\varphi(x)-\varphi(y)|\le\omega(\Omega_1,\,|x-y|),


\quad x,y\in\Delta.


\end{gather*}


\end{thm}





Заметим, что теорема 2, как видно из построений, остается справедливой,


если в ней вместо отрезка $\Delta$ взять любой компакт из $\Delta$.





Чтобы оценить $E_n(f,p,\Delta)$ для заданных на $\Delta=[a,b]$ функции


$f(x)$ и веса $p=(p_-,p_+)$, определим при


$\varepsilon>0$, как и выше, полунепрерывную сверху на множестве


$\Delta_+$ функцию $M_+(x)=M_+(x,\varepsilon)$ и


полунепрерывную снизу на множестве $\Delta_-$ функцию


$m_-(x)=m_-(x,\varepsilon)$ и для $\delta\ge 0$ рассмотрим величину


$$


\omega_\varepsilon(f,p,\delta)=\sup[M_+(x)-m_-(y)]^+,


$$


где супремум берется по $x\in\Delta_+$, $y\in\Delta_-$,


$|x-y|\le\delta$.


Легко видеть, что $\omega_\varepsilon(f,p,\delta)$ не убывает и


полунепрерывна сверху для $\delta\ge 0$;


непрерывна в точке $\delta=0$ и $\omega_\varepsilon(f,p,0)=0$, если


$M_+(x)\le m_-(x)$ при $x\in\Delta_\varepsilon\ne\emptyset$; при


$p_-(x)\equiv p_+(x)\equiv 1$ и непрерывной на


отрезке $\Delta$ функции $f(x)$ для всех $\varepsilon\in(0,1]$ величина


$\omega_\varepsilon(f,p,\delta)$ совпадает с


обычным равномерным модулем непрерывности $\omega(f,\delta)$ функции


$f(x)$ на $\Delta$.





Если $\dist(\Delta_-,\Delta_+)=d>0$, то доопределим величину


$\omega_\varepsilon(f,p,\delta)$ нулем для $\delta\in[0,d)$.





Для $\varepsilon>0$ рассмотрим семейство


$\omega=\{\omega_\varepsilon(\delta)\}$ функций типа модуля


непрерывности,


именно, при каждом $\varepsilon>0$ функция $\omega_\varepsilon(\delta)$


непрерывна и не убывает на $[0,+\infty)$, $\omega_\varepsilon(0)=0$,


$\omega_\varepsilon(\delta+h)\le\omega_\varepsilon(\delta)+


\omega_\varepsilon(h)$ ($\delta,h\ge 0$).


Далее при $x\in\Delta$ функции


\begin{align*}


\varphi_+(x,f)&=\sup\{M_+(y)-\omega_\varepsilon(|x-y|):


y\in\Delta_+\},\\


\varphi_-(x,f)&=\inf\{m_-(y)+\omega_\varepsilon|(x-y|):


y\in\Delta_-\}


\end{align*}


играют роль операторов ``сглаживания'' функций $f(x)$ ($x\in\Delta$) в


$p$-норме $|\cdot|_{p,\Delta}$.





Ниже для чисел $A$ и $\varepsilon>0$ через


$J(\varepsilon)=\{(\alpha,\beta)\}$ обозначим множество всех


дополнительных интервалов компакта $\Delta_\varepsilon$ и интервалов


$(a,\min\Delta_\varepsilon)$, $(\max\Delta_\varepsilon,b)$ и положим


\begin{align*}


& E_\pm(A)=\Delta(p_\pm\ge\varepsilon)\cap


\Delta(\pm f>\pm A+\varepsilon),\\


& E(A)=E_-(A)\cup E_+(A).


\end{align*}


Для непрерывной на $\Delta_\varepsilon$ функции $\varphi(x)$ будем


пользоваться также более кратким обозначением


$$


E_\gamma=E(\varphi(\gamma))\cap\,(\alpha,\beta);\ \;


\gamma=\alpha,\beta;\ \; (\alpha,\beta)\in J(\varepsilon).


$$


Для ограниченного на $\Delta=[a,b]$ веса $p=(p_-,p_+)$


по заданному числу


$L>0$ и семейству $\omega=\{\omega_\varepsilon(\delta)\}$ с


$\omega_\varepsilon(b-a)\le 2L$ определим класс


$K(L,p,\omega)$ функций $f(x)$ ($x\in\Delta$) с


$\|f\|_\Delta\le L$


таких,


что при $\delta>0$ и всех достаточно малых $\varepsilon>0$ выполняется


неравенство


$\omega_\varepsilon(f,p,\delta)\le\omega_\varepsilon(\delta)$.





\begin{thm}


Если на отрезке $\Delta=[a,b]$ функция $f\in K(L,p,\omega)$, то при


$n=1,2,\dotsc$ имеем


$$


E_n(f,p,\Delta)\le 6(L+\|p\|_\Delta)\inf_{\varepsilon>0}


\Bigl\{\varepsilon+\omega_\varepsilon\Bigl(


\frac{b-a}{n}\Bigr)\Bigr\}.


$$


\end{thm}





{\bf Доказательство теоремы 1. Необходимость.} Допустим, что существует


такое достаточно малое $\varepsilon>0$, что в некоторой точке


$x_0\in\Delta_\varepsilon$ имеем


$M_+(x_0)=M_+(x_0,\varepsilon)>m_-(x_0)=m_-(x_0,\varepsilon)$. Пусть


$\gamma>0$ такое, при котором


\begin{equation}


M_+(x_0)-3\gamma>m_-(x_0)+3\gamma.


\end{equation}


По определению $M_+(x_0)$ и $m_-(x_0)$ найдется такое $\delta_0>0$, что


при любом $\delta\in(0,\delta_0)$ существуют


точки $x$ и $y$ из интервала $(x_0-\delta,x_0+\delta)$, для которых


$p_+(x)\ge\varepsilon$,


$p_-(y)\ge\varepsilon$,


$M_+(x_0)<f(x)+2\gamma$, $m_-(x_0)>f(y)-2\gamma$, а следовательно, с


учетом (2) выполнено неравенство


\begin{equation}


f(x)-f(y)>2\gamma.


\end{equation}





С другой стороны, пусть $Q(x)$ --- полином, для которого


$|f-Q|_{p,\Delta}<\omega$, $0<\omega<\gamma\varepsilon/4$.


Выберем $\delta>0$ из условий


$|Q(t)-Q(\tau)|<\gamma/2$ при $|t-\tau|<2\delta$, $t,\tau\in\Delta$;


$\delta\in(0,\delta_0)$.


Тогда для точек $x$ и $y$ из (3) имеем


$$


f(x)-f(y)=[f(x)-f(y)]^+\le [f(x)-Q(x)]^+ +


[f(y)-Q(y)]^- +[Q(x)-Q(y)]^+ <\gamma


$$


в противоречие с (3).


\medskip





{\bf Достаточность.} Пусть $\varepsilon>0$ и выполнено условие (1).


Построим


вспомогательную функцию $F(x)$, непрерывную на $\Delta$. Так как


$M_+(x)=M_+(x,\varepsilon)$


и $m_-(x)=m_-(x,\varepsilon)$ полунепрерывны соответственно сверху и


снизу на $\Delta_\varepsilon$ и выполнено


неравенство (1), то по теореме об отделении полунепрерывных функций


непрерывными (\cite{7}, с.\,390), очевидно, существует непрерывная на


компакте $\Delta_\varepsilon$ функция


$\varphi(x)=\varphi(x,\varepsilon)$,


для которой $M_+(x)\le\varphi(x)\le m_-(x)$, $x\in\Delta_\varepsilon$.


Положим $F(x)=\varphi(x)$ при $x\in\Delta_\varepsilon$.





Из определения функций $M_+(x)$ и $m_-(x)$ следует, что принадлежащий


$\Delta_\varepsilon$ конец $\gamma$


интервала $(\alpha,\beta)\in J(\varepsilon)$ не может служить предельной


точкой множества $E(A)\cap\,(\alpha,\beta)$ ни при


каком $A\in[M_+(\gamma),\,m_-(\gamma)]$. Пусть $(\alpha,\beta)$ ---


произвольный дополнительный


интервал $\Delta_\varepsilon$. Так как


$M_+(\alpha)\le\varphi(\alpha)\le m_-(\alpha)$, найдется точка


$\alpha_1\in(\alpha,\beta)$ такая, что


$(\alpha,\alpha_1]\cap E_\alpha=\varnothing$.





Аналогичная точка $\beta_1\in(\alpha_1,\beta)$ найдется и для точки


$\beta$. При


$M_1=\max\{M(\alpha,\beta),\varphi(\alpha),\varphi(\beta)\}$,


$M_2=\min\{m(\alpha,\beta),\varphi(\alpha),\varphi(\beta)\}$ на


$[\alpha,\alpha_1]$ в


качестве графика функции $F(x)$ возьмем прямую, соединяющую точку


$(\alpha,\varphi(\alpha))$ с точкой $(\alpha_1,M_1)$ при


$[\alpha,\alpha_1]\subset\Delta_+$ и с точкой $(\alpha_1,M_2)$ при


$[\alpha,\alpha_1]\subset\Delta_-$.





Пусть теперь


$[\alpha,\alpha_1]\setminus(\Delta_+\cup\Delta_-)\ne\varnothing$. Тогда,


т.\,к. $(\alpha,\beta)\cap\Delta_\varepsilon=\varnothing$, найдется такой


интервал $(u,v)\subset(\alpha,\alpha_1)$, что $p_+(x)<\varepsilon$ и


$p_-(x)<\varepsilon$ при $x\in(u,v)$. Исправим точку $\alpha_1$, положив


$\alpha_1=v$, и рассмотрим точки


$\alpha_\pm=\inf\{\Delta_\pm\cap\,[\alpha_1,\beta_1]\}$.


Возьмем


$F(x)=\varphi(\alpha)$ при


$x\in(\alpha,u]\cap\,(\Delta_+\cup\Delta_-)$; $F(v)=M_1$, если


$\alpha_+<\alpha_-$ или точки $\alpha_-$ нет; $F(v)=M_2$, если


$\alpha_-<\alpha_+$ или $\alpha_-$ существует,


но $\alpha_+$ нет; затем продолжим $F(x)$ непрерывно на весь


$[\alpha,\alpha_1]=[\alpha,v]$, полагая ee линейной на оставшихся


интервалах.





На $[\beta_1,\beta]$ функция $F(x)$ определяется аналогично случаю


отрезка $[\alpha,\alpha_1]$.





На интервале $(\alpha_1,\beta_1)$ график функции $F(x)$ совпадает с


прямой, соединяющей точки\linebreak


$(\alpha_1,F(\alpha_1))$ и $(\beta_1,F(\beta_1))$, если


$(\alpha_1,\beta_1)\cap\,(\Delta_+\cup\Delta_-)=\varnothing$; в противном


случае положим $F(x)=M_1$ при $x\in(\alpha_1,\beta_1)\cap\Delta_+$,


$F(x)=M_2$ при $x\in(\alpha_1,\beta_1)\cap\Delta_-$ и


продолжим непрерывно на $[\alpha_1,\beta_1]$, считая ee линейной на


оставшихся интервалах.





На промежутках $[a,\min\Delta_\varepsilon)$ и


$(\max\Delta_\varepsilon,b]$ функция $F(x)$ строится вполне аналогично


рассмотренному случаю дополнительных интервалов множества


$\Delta_\varepsilon$.





Для построенной непрерывной на $\Delta$ функции $F(x)$ по теореме


Вейерштрасса существует алгебраический полином $Q(x)$ такой, что


$|F(x)-Q(x)|<\varepsilon$, $x\in\Delta$.





Возьмем в качестве $Q(x)$ полином (соответствующей степени) наилучшего


равномерного приближения для $F(x)$ на $\Delta$ и будем считать


$\varepsilon\le\|f\|_\Delta$. Тогда


$\|Q\|_\Delta\le\varepsilon+\|F\|_\Delta\le 2\|f\|_\Delta$.





Оценим сверху величину $r(x)=|(f-Q,p)(x)|$, $x\in\Delta$.





Для $x\in\Delta_\varepsilon$ рассмотрим четыре случая:


\begin{itemize}


\item[1)] если $x\in F_\varepsilon$, то $f(x)\le M_+(x)\le\varphi(x)\le


m_-(x)\le f(x)$, поэтому


$r(x)=|(F-Q,p)(x)|\le\|p\|_\Delta\varepsilon$;


\item[2)] если $p_+(x)<\varepsilon$ и $p_-(x)<\varepsilon$, то


$r(x)\le 3\|f\|_\Delta\varepsilon$;


\item[3)] если $p_+(x)\ge\varepsilon$, а $p_-(x)<\varepsilon$, то


$f(x)\le M_+(x)\le\varphi(x)=F(x)$, поэтому


$r(x)=[Q(x)-f(x)]p_-(x)\le 3\|f\|_\Delta\varepsilon$ при


$f(x)\le Q(x)$, а при $f(x)>Q(x)$


имеем


$r(x)=[f(x)-Q(x)]p_+(x)\le[F(x)-Q(x)]p_+(x)


\le\|p\|_\Delta\varepsilon$;


\item[4)] если $p_+(x)<\varepsilon$, $p_-(x)\ge\varepsilon$, то


поступаем аналогично предыдущему случаю.


\end{itemize}





Пусть теперь $x$ --- точка дополнительного интервала $(\alpha,\beta)$


множества $\Delta_\varepsilon$ и рассмотрим возможные различные случаи.





Пусть $x\in[\alpha,\alpha_1]\subset\Delta_+$. Если


$p_+(x)<\varepsilon$ и $p_-(x)<\varepsilon$, то


$r(x)\le 3\|f\|_\Delta\varepsilon$; если же


$p_+(x)\ge\varepsilon$, $p_-(x)<\varepsilon$, то


$f(x)\le\varphi(\alpha)+\varepsilon$, поэтому при


$Q(x)\ge f(x)$ и $Q(x)<f(x)$ соответственно имеем


$r(x)=[Q(x)-f(x)]p_-(x)\le 3\|f\|_\Delta\varepsilon$,


$r(x)=[f(x)-Q(x)]p_+(x)\le [\varphi(\alpha)+\varepsilon-Q(x)]


\|p\|_\Delta\le [F(x)+\varepsilon-Q(x)]\|p\|_\Delta\le


2\varepsilon\|p\|_\Delta$ (по выбору


точки $\alpha_1$ другие случаи невозможны).





Случай $x\in[\alpha,\alpha_1]\subset\Delta_-$ аналогичен предыдущему.





Рассмотрим случай


$[\alpha,\alpha_1]\setminus(\Delta_+\cup\Delta_-)\ne\varnothing$. Пусть


$(u,v)$ --- интервал, о котором шла речь выше при


построении $F(x)$ для этого случая, причем считаем $v=\alpha_1$, и


возьмем $x\in(\alpha,v)$.


Тогда возможны три варианта: 1)~$p_+(x)<\varepsilon$,


$p_-(x)<\varepsilon$;


2)~$p_+(x)\ge\varepsilon$, $p_-(x)<\varepsilon$;


3)~$p_+(x)<\varepsilon$,


$p_-(x)\ge\varepsilon$. В этих вариантах, как и выше, либо


$r(x)\le 3\|f\|_\Delta\varepsilon$, либо


$r(x)\le 2\|p\|_\Delta\varepsilon$.





Оценки $r(x)$ при $x\in(\beta_1,\beta)$ вполне аналогичны оценкам при


$x\in(\alpha,\alpha_1)$.





Для $x\in[\alpha_1,\beta_1]$ рассмотрим три возможных случая:


\begin{itemize}


\item[1)] $p_+(x)<\varepsilon$, $p_-(x)<\varepsilon$;


\item[2)] $p_+(x)\ge\varepsilon$, $p_-(x)<\varepsilon$;


\item[3)] $p_+(x)<\varepsilon$, $p_-(x)\ge\varepsilon$.


\end{itemize}


В результате получим $r(x)\le 3\|f\|_\Delta\varepsilon$ или


$r(x)\le\|p\|_\Delta\varepsilon$.





В случае


$x\in[a,\min\Delta_\varepsilon)\cup\,(\max\Delta_\varepsilon,b]$ оценки


$r(x)$ аналогичны вышеприведенным.





Следовательно, если $\Delta_\varepsilon\ne\varnothing$, то при


$x\in\Delta$ имеем


$$


|(f-Q,p)(x)|\le\max\{3\|f\|_\Delta,\,2\|p\|_\Delta\}\varepsilon.


$$


Если же $\Delta_\varepsilon=\varnothing$, то требуемое легко получим,


взяв кусочно-линейное непрерывное продолжение на $\Delta$ такой


функции $F(x)$,


что


\begin{align*}


F(x)&=\sup\{f(t):t\in\Delta(p_+\ge\varepsilon)\},\quad x\in\Delta_+;\\


F(x)&=\inf\{f(t):t\in\Delta(p_-\ge\varepsilon)\},\quad x\in \Delta_-.


\end{align*}


Достаточность доказана. $\quad\square$


\medskip





{\bf Доказательство теоремы 2.} Пусть, например,


$\varphi(x)=\varphi_u(x)$. При $x,y\in\Delta$ по определению и


свойствам $\Omega(\delta)$ имеем


\begin{gather*}


u(x)-v(y)\le[u(x)-v(y)]^+\le \Omega(|x-y|)\le


\Omega_1(|x-y|)\le \omega(\Omega_1,\,|x-y|);\\


%


u(x)-\omega(\Omega_1,\,|x-y|)\le v(y).


\end{gather*}





Переходя к супремуму при $x\in\Delta$, получим $\varphi(y)\le v(y)$.


Неравенство $u(y)\le\varphi(y)$


($y\in\Delta$) получается при $y=x$ непосредственно из определения


$\varphi_u(x)$.





Ввиду полунепрерывности сверху функции $u(x)$ для каждой точки


$x\in\Delta$ существует точка $z(x)\in\Delta$ такая, что


$$


\varphi(x)=\varphi_u(x)=u(z(x))-\omega


(\Omega_1,\,|x-z(x)|).


$$


Тогда при $x,y\in\Delta$ получим


\begin{multline*}


\varphi(x)\ge u(z(y))-\omega(\Omega_1,\,|x-z(y)|)\ge \\


\ge u(z(y))-\omega(\Omega_1,|x-y|)-


\omega(\Omega_1,\,|y-z(y)|)= \varphi(y)-


\omega(\Omega_1,\,|x-y|).


\end{multline*}


Следовательно, $\varphi(y)-\varphi(x)\le\omega(\Omega_1,\,|x-y|)$.


Остается поменять местами $x$ и $y$, чтобы получить второе из требуемых


неравенств. $\quad\square$


\medskip





{\bf Доказательство теоремы 3} проводится по аналогии с доказательством


достаточности в теореме 1 и с доказательством теоремы 2.





Отметим лишь, что для $f\in K(L,p,\omega)$ при каждом достаточно малом


$\varepsilon>0$ функция


$F(x)$, равная на $\Delta$ любой из функций $\varphi_+(x,f)$ и


$\varphi_-(x,f)$, построенных выше, обладает свойствами


\begin{gather*}


|F(x)-F(y)|\le\omega_\varepsilon(|x-y|),\quad x,y\in\Delta;\\


M_+(x)\le f(x),\quad x\in\Delta_+;\\


m_-(x)\ge F(x),\quad x\in\Delta_-;\\


|f-F|_{p,\Delta}\le 4L\varepsilon.\quad\square


\end{gather*}





В заключение отметим, что в \cite{6} и \cite{8} рассмотрен


случай непрерывного знакочувствительного веса.





\begin{thebibliography}{9}





\bibitem{1}


Долженко Е.П., Севастьянов Е.А. {\em Знакочувствительные аппроксимации.


Пространство знакочувствительных весов. Жесткость и свобода системы} //


Докл. РАН. -- 1993. - Т.\,332. -- \No\,6. -- С.\,686--689.


\bibitem{2}


Долженко Е.П., Севастьянов Е.А. {\em Знакочувствительные аппроксимации.


Вопросы единственности и устойчивости} // Докл. РАН. -- 1993. -- Т.\,333.


-- \No\,1. -- С.\,5--7.


\bibitem{3}


Бабенко В.Ф. {\em Несимметричные приближения в пространствах


суммируемых функций} // Укр. матем. журн. -- 1982. -- Т.\,34. -- \No\,4.


-- С.\,409--416.


\bibitem{4}


Крейн М.Г., Нудельман А.А. {\em Проблема моментов Маркова и


экстремальные задачи. Идеи и проблемы П.Л.\,Чебышева и А.А.\,Маркова и


их дальнейшее развитие}. -- М.: Наука, 1973. -- 551~с.


\bibitem{5}


Симонова И.Э., Симонов Б.В. {\em О полиноме наилучшего несимметричного


приближения в пространстве Орлича} // Изв. вузов. Математика. -- 1993. --


\No\,11. -- С.\,50--56.


\bibitem{6}


Рамазанов А.-Р.\,К. {\em О прямых и обратных теоремах теории


аппроксимации в метрике знакочувствительного веса} // Anal. math. --


1995. -- V.\,21. -- P.\,191-212.


\bibitem{7}


Натансон И.П. {\em Теория функций вещественной переменной}. -- 3-е изд.


-- М.: Наука, 1974. -- 480~с.


\bibitem{8}


Рамазанов А.-Р.\,К. {\em Рациональная аппроксимация со знакочувствительным


весом} // Матем. заметки. -- 1996. -- Т.\,60. -- \No\,5. --


C.\,715--725.








\end{thebibliography}





\vskip 2cm





\post{\it Дагестанский государственный}{\it Поступила}


\post{\it университет}{08.07.1996}





\label{end}


\end{document}


